% §6.3 LEACE + IGBP deconfounding of BERT embeddings, 12-city panel.
% Numbers locked against results/leace12_igbp/{city}_leace.json + leace_12city_table.csv
% (full-corpus continuous-latlon LEACE+IGBP rerun, float32 GPU). Supersedes v2,
% which was locked against the stale small-corpus results/leace_12city/.

\subsection{Spatial deconfounding of BERT embeddings via LEACE and IGBP}
\label{sec:leace-12}

The sentence-BERT representation that underlies the Baur channel in
Section~\ref{sec:baur-12city} encodes more than the language a listing agent
uses. Pre-trained on a generic corpus, the encoder has no particular incentive
to separate semantic content about a property's condition or amenities from the
geographic information that correlates with its neighborhood. If the embedding
carries location-targeted information into the leading principal component, the
per-standard-deviation effect reported in Section~\ref{sec:baur-12city} inherits
a spatial confounder mediated through the encoder, and the identification claim
is correspondingly weaker. We assess how far this concern can be addressed by
erasing the lat-lon concept from the representation, and we report honestly where
the erasure is complete and where it is only partial.

The deconfounding runs a linear projection and then a nonlinear correction, and
the linear step carries most of the load. Applying the LEACE projection of
\citet{belrose2023leace} to the BERT embedding removes the linear subspace that
the mean-conditional-covariance estimator identifies with the concept variable,
here the latitude-longitude coordinate pair of the listing, and it comes with the
exact-guardedness guarantee that no linear predictor of the residual embedding can
recover the concept any better than a constant baseline would. That guarantee
holds in every market we examine, with the out-of-fold ridge $R^2$ for predicting
the coordinate from the embedding falling from a raw value that reaches $0.50$ in
New York and $0.42$ in Dallas to within $0.002$ of zero after erasure, and the
linear-guardedness residual averaging $2.7\text{e-}8$ across the twelve markets,
ranging from $8.3\text{e-}9$ to $6.0\text{e-}8$, which is the floating-point floor
of the single-precision computation, so the linear erasure comes out complete
throughout the panel.

The concept is routed through the continuous coordinate representation rather
than the categorical ZIP-code identifier in every market. Categorical LEACE has
known instability when the categorical cardinality is large, and our panel ranges
from $30$ ZIP codes in San Francisco to $158$ in New York. Although the
per-market test folds are now large, from $296$ listings in San Francisco to
$4{,}569$ in New York, the continuous routing sidesteps the cardinality question
entirely without discarding the spatial information the projection is designed to
remove, since the coordinate pair is a sufficient statistic for location at the
resolution that matters here.

What a linear projection cannot reach is the location the embedding encodes
nonlinearly, and the iterative gradient-based correction of
\citet{iskander2023shielded} is aimed precisely there, alternating between fitting
a multilayer-perceptron probe of the coordinate on the residual and projecting out
the direction in which that probe gathers its predictive mass. The correction cuts
the nonlinear leakage substantially without removing it, taking the average
nonlinear probe $R^2$ from $0.44$ on the raw embedding down to $0.20$ after five
outer passes, a reduction of roughly a half, while the residual that remains
varies sharply across markets. San Francisco is left with a nonlinear $R^2$ of
$0.01$, which is erasure for any practical purpose, whereas New York, Dallas, and
Phoenix hold onto residuals near $0.30$ to $0.34$ and Portland near $0.26$, and
the markets that keep the most are the larger corpora whose listing language
carries the richest location-targeted structure, which is also where a nonlinear
probe has the most left to find.

\paragraph{What the diagnostic does and does not settle.} The Baur estimates of
Table~\ref{tab:baur-12city} are computed on the unprojected embedding, in keeping
with the published Baur, Rosenfelder, and Lutz protocol, so the erasure here is a
diagnostic of what the representation contains rather than a change to any
per-market number. As a diagnostic it speaks cleanly to the part of the problem
the linear treatment can actually express, since the leading principal component
is a linear functional of the embedding and the linear erasure is exact in every
market, so whatever location a linear treatment could carry has been shown to be
removable in full, while the nonlinear residual that survives IGBP in the larger
markets bounds the rest, fixing how much location a nonlinear function of the
deconfounded representation could still smuggle back. Showing that the embedding
encodes geography is not, however, the same as showing that the price effect is
confounded by it, and the diagnostic on its own cannot close that gap, because a
representation may carry a confounder along directions that have nothing to do
with the one on which it acts. The question of whether the location the embedding
demonstrably encodes is the location that actually moves the price is taken up
directly in Section~\ref{sec:decomposition}, where the text treatment is held
fixed and geography is moved on and off the controls around it.
